\textbf{Время:} 01:45

\textbf{Тезис:} независимое усреднение двух точек приводит к квадрату модуля Фурье-образа источника.

\textbf{Визуал:} повторять выбор A и B → двойное среднее → факторизация.

\textbf{Текст:} Плотность rho описывает вероятности точек и времён испускания при выбранном среднем импульсе K. Её интеграл равен единице.

Усредняем косинус по всем парам точек. При независимом испускании вероятность пары — произведение двух плотностей, поэтому возникает двойной интеграл. Экспонента разности фаз раскладывается на множитель для A и множитель для B. Интегралы оказываются сопряжёнными; произведение равно квадрату модуля одного интеграла. Так появляется преобразование Фурье.

Нужны хаотическое испускание, плоские волны и гладкая импульсная зависимость источника, позволяющая использовать общий K. При малом q фазы мало различаются; при большом взаимно усредняются. Большим расстояниям достаточно меньшего изменения q для такой компенсации.

Пауза: модуль теряет фазу, а энергия связана с импульсом. Поэтому произвольное четырёхмерное изображение источника однозначно не восстанавливается.

\textbf{Переход:} Проверим смысл этой связи на источнике с простой известной формой.
